\documentclass[10pt,a4paper]{amsart}
\usepackage[margin=19mm]{geometry}
\usepackage{amsmath,amssymb,mathtools}
\usepackage[scr=rsfs,cal=euler]{mathalfa}
\usepackage{xfrac}
\usepackage[unicode,hidelinks,bookmarks=false]{hyperref}
\newcommand{\ie}{i.e.\ }
\newcommand{\cf}{cf.\ }
\newcommand{\resp}{respectively\ }
\newcommand{\Subsection}{\subsection}
\providecommand{\nobreakdash}{\nobreak\hbox{-}}
\setlength{\emergencystretch}{2em}
\multlinegap0pt
\usepackage{bm}
\usepackage{bbm}
\usepackage{dsfont}
\usepackage{xspace}
\usepackage{cancel}
\usepackage{mathtools}
\usepackage{amsthm}
\makeatletter
\@ifundefined{theorem}{\newtheorem{theorem}{Theorem}}{}
\@ifundefined{lemma}{\newtheorem{lemma}[theorem]{Lemma}}{}
\@ifundefined{question}{\newtheorem{question}[theorem]{Question}}{}
\makeatother
\def\N{{\mathbb N}}
\def\Z{{\mathbb Z}}
\title[A refined structure theorem for polynomial return-time sets ]{A refined structure theorem for polynomial return-time sets in minimal systems}
\author{Ioannis Kousek}
\date{Source: arXiv:2607.18987v2 (CC BY 4.0)}
\begin{document}
\maketitle

\noindent\emph{Source and licence.} A refined structure theorem for polynomial return-time sets in minimal systems. Version arXiv:2607.18987v2 (CC BY 4.0), distributed under the Creative Commons Attribution 4.0 International licence, as recorded in the arXiv licence metadata for this exact version and shown on the abstract page https://arxiv.org/abs/2607.18987v2. Full rights notice: licenses/arxiv-2607.18987-CC-BY-4.0.txt.\par\smallskip

\noindent\textit{Editorial scope.} Counted unit: the Question whose source label is \texttt{Qiu strengthened} in the article (source0.tex lines 1115--1119, source label \texttt{Qiu strengthened}), delivered together with the article's own complete original proof of it (source0.tex lines 1123--1148, one \texttt{proof} environment of 26 lines). No mathematics is written, repaired or re-derived: statement and proof are the article's own text line for line. The proof is detached from the statement in the source: the article states the result at lines 1115--1119 and prints its proof at lines 1123--1148, and the proof environment's own header names the statement, so the association is the article's own. Required context retained uncounted: A-B pws (lines 427--429); GKMLRR question (lines 1046--1056). Every definition, display and label of those lines is retained unchanged. Omitted from this package and not redistributed: 1--426, 430--1045, 1057--1114, 1120--1122, 1149--1394. Nothing inside a retained excerpt is omitted or altered: every retained body is delivered exactly as the article's lines are written, so no omission receipt of any kind accompanies this package. Citations: the retained excerpts carry no citation command at all, so no bibliography entry of the article is transcribed and no reference is redistributed. Reference adaptation: none was needed and none was made; every reference inside the delivered unit resolves inside it, so no unresolved reference or citation is produced. Printed slips kept literal: no printed word, symbol, hypothesis or number of the counted statement or of its proof was altered. Layout and class: the article's own class is \texttt{amsart} with packages including scrextend, fontenc, pgf, pgfarrows, pgfnodes, pgfautomata, pgfheaps, pgfshade, hyperref, amssymb, tikz-cd, amsmath, babel, cleveref; the wrapper is the lane's pinned \texttt{amsart} editorial wrapper at 10pt on A4, which restates the theorem-like environments the retained text needs and adds the article's own macro definitions that the retained text needs (\texttt{\textbackslash N}, \texttt{\textbackslash Z}), with extra wrapper packages (bm, bbm, dsfont, xspace, cancel, mathtools). The delivered numbering is the unit's own. Assets: no retained span contains a figure environment, a graphics-inclusion command, a PDF-page inclusion, a table or a TikZ picture; no image file is copied into this package, the package declares no asset dependency and no image file is redistributed with it. Licence: version arXiv:2607.18987v2 of this article is distributed under the Creative Commons Attribution 4.0 International licence, as recorded in the arXiv licence metadata for this exact version and shown on the abstract page https://arxiv.org/abs/2607.18987v2. A full rights notice is delivered at licenses/arxiv-2607.18987-CC-BY-4.0.txt. Characters: every retained excerpt is pure ASCII, so the delivered engine, XeLaTeX with its default Latin Modern text font, reports no missing character.
\begin{lemma}\label{A-B pws}
Let $B\subset A\subset \Z$. If the set $\bigcap_{n\in F}B-f$ is syndetic for any finite $F\subset A$, then $A\setminus B$ is not piecewise syndetic. 
\end{lemma}

\begin{question}\label{GKMLRR question}
Let $(X,T)$ be a minimal and invertible
topological dynamical system. Denote by $(X_{\infty},T)$ its 
maximal infinite-step pronilfactor, and let $\pi: X \to X_{\infty}$ be the associated factor map. Is it true that for all
$d\in \N$, all nonempty, open $U_1,\ldots,U_d \subset X$, and all 
essentially distinct $p\in \Z[x]^d$, the set 
\begin{equation}\label{set difference of return times question}
R_p(\pi U_1,\ldots,\pi U_d) \setminus R_p(U_1,\ldots,U_d)     
\end{equation}   
is not piecewise syndetic?
\end{question}

\begin{question}\label{Qiu strengthened}
Let $(X,T)$ be a minimal and invertible
topological dynamical system. Denote by $(X_{\infty},T)$ its maximal infinite-step pronilfactor, and let $\pi: X \to X_{\infty}$ be the associated factor map. For all $d\in \N$, all nonempty, open $U_1,\ldots,U_d \subset X$, and all 
essentially distinct $p\in \Z[x]^d$, if $R_p(\pi U_1,\ldots,\pi U_d )$ is piecewise syndetic, is it true that $R_p(U_1,\ldots,U_d) \neq \emptyset$? 
\end{question}

\begin{proof}[Proof that positive answer to Question \ref{Qiu strengthened} implies positive answer to Question \ref{GKMLRR question}.]
\ \\
Let $A=R_p(\pi U_1,\ldots,\pi U_d )$ and $B=R_p(U_1,\ldots,U_d)$ 
and note that $B\subset A$. Hence, if $A$ is not piecewise 
syndetic, the implication is trivial, which is why we restrict to 
this assumption.
In view of Lemma \ref{A-B pws}, it suffices to show that 
$\bigcap_{a\in F} (B-a)$ is syndetic, for any $F\subset A$ finite. 

Take $a\in A$. Then, 
$$\pi(T^{-p_1(a)}U_1) \cap \cdots \cap \pi(T^{-p_d(a)}U_d) \neq \emptyset,$$
and by the conditional result applied to the polynomials $n\mapsto p_i(n+a)-p_i(a)$, $i=1,\ldots,d$, and the open sets $V_i=T^{-p_i(a)}U_i$, it follows that there exists $n\in \N$ such that
$$T^{-p_1(n+a)}U_1 \cap \cdots \cap T^{-p_d(n+a)}U_d \neq \emptyset. $$

Now, take any finite collection $a_1,\ldots,a_F \in A$. Then, by 
minimality, we can find
$\ell_1,\ldots,\ell_{F}\in \N$ such that 
$$\bigcap_{j=1}^{F} \bigcap_{i=1}^{d}  \pi(T^{-(p_i(a_j)+ \ell_j)}U_i) \neq \emptyset$$
and again, by the conditional result applied to the polynomials $n\mapsto p_i(n+a_j)-p_i(a_j)$, and the open sets $V_{i,j}=T^{-(p_i(a_j)+ \ell_j)}U_i$, for $i=1,\ldots,d$ and $j=1,\ldots,F$, it follows that there exists $n\in \N$ such that
$$\bigcap_{j=1}^{F} \bigcap_{i=1}^{d}  T^{-(p_i(n+a_j)+ \ell_j)}U_i) \neq \emptyset.$$ 
In particular, by minimality, the set of $n\in \N$ for which the latter is nonempty is syndetic. Finally, note that any such $n$ is contained in $B-a_j$, for any $j=1,\ldots,F$, so the result follows. Indeed, 
$$\bigcap_{i=1}^{d}  T^{-(p_i(n+a_j)+ \ell_j)}U_i) \neq \emptyset$$
means that there exists $x\in X$ such that 
$$T^{\ell_j}x \in T^{-p_1(n+a_j)}U_1 \cap \cdots \cap T^{p_d(n+a_j)}U_d,$$
and hence $n+a_j\in R_p(U_1,\ldots,U_d)=B$.
\end{proof}

\end{document}
